\section{Simulation study}\label{sec:design}

Every scenario, comparison, threshold and decision rule below was written down and fixed, with a
cryptographic hash, before any of its results were computed. The specifications, their hashes and
every per-replicate $p$-value are in the reproducibility archive \citep{edge2_repro_2026}; the
comparisons that were specified in advance, and their outcomes, are listed in Supporting Information
S0. Section~\ref{sec:reading} lists where the analysis departs from them.

\subsection{Designs and departures}\label{sec:families}

The working model is a logistic regression for a binary outcome with one to three covariates. Data are
generated from a truth that departs from it in one of six families:
\begin{enumerate}
\item \textbf{Symmetric tails}: the cauchit and $t_4$ latent links, and the symmetric members of
      Stukel's family, heavier or lighter than the logistic at both ends.
\item \textbf{Asymmetric links}: the complementary log--log and log--log links and the asymmetric
      members of Stukel's family, which bend the curve once.
\item \textbf{The Stukel symmetric family} at its published shape parameters.
\item \textbf{Omitted terms}: a quadratic term, an interaction, or a covariate left out of the model.
\item \textbf{Rough misfit}: an oscillating calibration curve, which no smooth basis can follow.
\item \textbf{Off-index departures}, which do not act through the linear predictor.
\end{enumerate}
Sample sizes run from $100$ to $50{,}000$. Each family is crossed with three regimes: a base design
(area under the curve about $0.77$), a high-discrimination design (about $0.89$), and a design with
$12\%$ events, the sparse regime in which partition tests are known to struggle. There are $131$ null
and $158$ alternative scenarios, with between $500$ replicates in the largest samples and $10{,}000$ in the
smallest.

\subsection{Comparators}\label{sec:comparators}

\begin{itemize}
\item The Hosmer--Lemeshow statistic and its Osius--Rojek/Farrington correction \EFname{}
      \citep{hosmer1980goodness,osius1992normal,farrington1996}, at $G=10$ and at the default
      partition; the Hosmer--Lemeshow statistic with equal-width risk groups
      \citep{hosmer1980goodness}; the Pigeon--Heyse test \citep{pigeon1999b}; and the tests of
      \citet{Tsiatis1980} and \citet{xie2008increasing}, which group the records into ten $K$-means
      clusters in covariate space, Tsiatis's through a score test on the cluster indicators. These four
      run at ten groups.
\item Stukel's score test in its joint and one-parameter forms, and its likelihood-ratio refit
      \citep{stukel1988generalized}. The form distributed in current software \citep{LogisticDx}
      refers the sum of two correlated score components to a chi-squared distribution on two degrees of
      freedom; that reference is not correct, with a median size of $0.066$ over the $154$ null settings in which it was run, those of the
      supplementary studies included
      (quartiles $0.049$ to $0.070$) against $0.049$ for the joint score, so the joint
      score form is used throughout.
\item The GiViTI calibration belt \citep{nattino2014new,nattino2015new}.
\item The cubic calibration likelihood-ratio test, which adds the square and cube of the fitted
      linear predictor to a logistic regression of the outcome on it and compares the fits on two
      degrees of freedom: the polynomial calibration model of the GiViTI belt with its degree fixed at
      three, which extends Pregibon's goodness-of-link test
      \citep{pregibon1980goodness,nattino2014new,nattino2015new}.
\item The kernel-smoothed residual test of \citet{le1991goodness,le1995goodness}, at its default kernel
      and bandwidth. We use the implementation in \texttt{ebrahim.gof}, because the implementation it
      was derived from omits a transpose that a weighted fit requires and is mildly liberal as a result
      \citep{ebrahim2026benchmark}.
\item The Liu projection test \citep{liu2024comprehensive}, a logistic form of the statistic of
      \citet{escanciano2006consistent}, with $B=250$ model-based bootstrap draws, from our implementation, which is in the archive and in
      \texttt{ebrahim.gof} from version 2.8.0.
\item BAGofT \citep{BAGofT2019} from its published package at its defaults: one hundred splits, one
      hundred simulations and the random-forest partition. One line of the package fails on models with a
      single covariate and was repaired; with two or more covariates the repaired function returns
      identical results.
\end{itemize}
The weighted grouped tests of \citet{hosmer2002goodness}, which multiply each record's residual by a
weight chosen for power against a specified alternative, are the closest earlier directed grouped
tests. Their weights are the residuals of a weighted regression of the alternative's direction on the
covariates, so they can grow with a record's covariates and Proposition~\ref{prop:bounded} does not
cover them; they are compared, with the paper's general-purpose weight, in the study of
Section~\ref{sec:whatprotects}. Every test in the list except the cubic calibration test is available in \texttt{ebrahim.gof} \citep{ebrahim2025gof}, where one call to
\texttt{run.all.gof()} fits the model once and returns all of them; \citet{ebrahim2026benchmark}
benchmark that battery under sparse data. The three tests outside the
partition family are slow, so they were run on subsets of the scenarios: le Cessie's test on $82$, the
projection test on $45$ and BAGofT on $32$. Every test sees the same data sets.

\subsection{Corrupted records}\label{sec:contamination-design}

In each replicate $k$ of the $n$ records are corrupted, with $k/n$ from $0.1\%$ to $1\%$ and $n$ equal
to $1000$ or $5000$:
\begin{description}
\item[C1, an exaggerated covariate.] The covariate of $k$ random records is multiplied by four, and
      their outcome is still drawn from the truth at the original value. The \emph{prediction} is wrong
      and the outcome is not.
\item[C1$'$, a sign error.] As C1, with the covariate multiplied by $-4$. The prediction is reversed
      rather than exaggerated.
\item[C2, missed events.] The $k$ records of highest true risk have their outcome set to zero. The
      prediction is right and the \emph{outcome} is wrong.
\item[C3, a unit error.] The covariate of $k$ records is divided by ten, which makes their
      predictions milder rather than extreme.
\item[C4, duplicated records.] The $k$ highest-risk records are copied over $k$ random rows, which
      checks that the grouped reference survives ties.
\end{description}
Formally, the null hypothesis is that the model is calibrated for the patients as they are,
$\Pr(y=1\mid x)=\pi(x;\beta)$ at the true covariates, while $k$ of the $n$ records are observed with
covariates $\tilde x\neq x$ and so receive the prediction $\pi(\tilde x;\hat\beta)$. Under a logistic truth
the model is correct for every uncorrupted record, so a rejection is a false alarm: the test is
reporting the corrupted records, not a miscalibration of the model for the population. It is not a type-I error in the strict sense, because a contaminated sample is a mixture; what
is measured is each test's stability of level over a small gross-error neighbourhood of the model
\citep{huber1964robust,hampel1986robust}. It is not the classical measurement-error problem, in which
every covariate is observed with error and the fit itself is biased \citep{carroll2006measurement};
here a few records carry a gross error and the fit moves by the amounts given in
Section~\ref{sec:mechanism}. Under a probit and a complementary log--log truth the same corruptions are
applied to a real misfit, to measure the cost in power. A value corrupted this way lies outside the
range of the clean covariate, and a range check on that variable would find it; the cohort of
Section~\ref{sec:realdata} shows an error that no range check finds. Wrong outcomes (C2) are treated differently: the recorded data are then miscalibrated, and a
calibration test should detect them. A further study compares EDGE with robust quasi-deviance versions
of Stukel's test \citep{cantoni2001robust}, and two others vary the partition with the data held fixed
and compare partitions that pool the extreme groups at contamination rates up to $4\%$. A last study, added after the main results were read and hashed before it ran, runs on the same data
sets EDGE, an ungrouped score test for the same
cubic directions in the predicted risk, EDGE with one record a group, Spiegelhalter's $z$
\citep{spiegelhalter1986probabilistic}, the weighted grouped tests of \citet{hosmer2002goodness}, and
each test after the records flagged by Pregibon's diagnostics \citep{pregibon1981logistic} are removed;
it adds misfit confined to the top $2\%$ of risk, to measure what protection costs, and a design with
twelve covariates at $n=2000$.

\subsection{Reading the results}\label{sec:reading}

Size is checked before power is reported. Power is size-adjusted: each test is read at the critical
value from its own matched null scenario, so every test is compared at its own realised level; where
a table reports raw rejection rates, it says so. Comparisons are paired: two tests are compared on the
same replicates by an exact McNemar test on the discordant pairs, with Holm's correction within each
family of comparisons. A test \emph{leads or ties} a rival in a scenario when its size-adjusted power
is no more than $0.01$ below the rival's; the specification allowed $0.014$, and the stricter margin is
used everywhere. A replicate on which a test returns no $p$-value counts as a non-rejection. The two
resampling tests return $p$-values on grids of $1/100$ and $1/250$ and reject at $p<\alpha$; every other
test rejects at $p\le\alpha$. Under corruption we say a test \emph{holds} while its false-alarm rate
stays below $0.10$, twice the nominal level; that is the bar the specification set. It is looser than
the band of three Monte Carlo standard errors used for size, and the tables give every rate so that
the reader can apply either.

Three readings depart from the specifications. The robust-fit study (Section~\ref{sec:robustest})
reads the contaminated columns of the tests whose level survived the robust fit, although its rule was
to read none once any test failed that check. The corruption experiment on the cohort was repeated on
a second covariate, chosen by a rule fixed before it was run, after the first showed no effect. And
the cohort analysis of Section~\ref{sec:realdata} keeps one admission per patient and splits by
patient, where the specification split admissions; the change was made because patients recur across
admissions.
